\clearpage\thispagestyle{empty}
\pdfbookmark[0]{Quellen zum Ceres- und Pallas-Bericht vom März 1804}{cerespallas180403-sources}
\begin{center}\textbf{Quellen zum Ceres- und Pallas-Bericht vom März 1804}\end{center}
\noindent Diese KI-Übertragung in heutiges Deutsch folgt dem Auszug aus dem Bericht von Franz Xaver von Zach in \href{https://gdz.sub.uni-goettingen.de/id/PPN236007467}{Gauß’ \textit{Werken}, Band~VI} (Göttingen, 1874), S.~246--247. Der Bericht beginnt auf \href{https://api.digitale-sammlungen.de/iiif/image/v2/bsb10538601_00260/full/1200,/0/default.jpg}{S.~246} als Beitrag XXV im Märzheft 1804 der \textit{Monatlichen Correspondenz}, Band~IX. Gauß’ Ephemeride für Pallas steht auf \href{https://api.digitale-sammlungen.de/iiif/image/v2/bsb10538601_00261/full/1450,/0/default.jpg}{S.~247}, \href{https://api.digitale-sammlungen.de/iiif/image/v2/bsb10538601_00262/full/1450,/0/default.jpg}{S.~248} und \href{https://api.digitale-sammlungen.de/iiif/image/v2/bsb10538601_00263/full/1450,/0/default.jpg}{S.~249}. Die \textit{Werke} beschreiben die Ephemeride in eckigen Klammern und geben die \href{https://api.digitale-sammlungen.de/iiif/image/v2/bsb10538601_00264/full/1450,/0/default.jpg}{Elemente und den Schlussabsatz auf S.~250} wieder. Die Liste der Sterne aus Piazzis Katalog beginnt in der Zeitschrift auf \href{https://api.digitale-sammlungen.de/iiif/image/v2/bsb10538601_00265/full/1250,/0/default.jpg}{S.~251}.
\par\smallskip\noindent Der Bericht verweist auf Gauß’ \href{https://api.digitale-sammlungen.de/iiif/image/v2/bsb10538600_00380/full/1450,/0/default.jpg}{Ceres-Ephemeride vom Oktober 1803}, die Ankündigung von Karl Ludwig Hardings \href{https://api.digitale-sammlungen.de/iiif/image/v2/bsb10538600_00547/full/1450,/0/default.jpg}{Karte des Laufs der Ceres} im Dezember und die zugehörigen \href{https://api.digitale-sammlungen.de/iiif/image/v2/bsb10538600_00548/full/1450,/0/default.jpg}{Vergleichssterne für Ceres}. Johann Elert Bodes Ephemeride für Pallas beginnt auf \href{https://api.digitale-sammlungen.de/iiif/image/v2/bsb10538335_00101/full/1450,/0/default.jpg}{S.~91 des \textit{Astronomischen Jahrbuchs für 1806}}.
\par\smallskip\noindent Übertragung und Satz wurden mit KI erstellt. Die Digitalisate der Zeitschrift und des Jahrbuchs stammen von der Bayerischen Staatsbibliothek.
\clearpage
